\chapter{Aplicación web de visualización}
\label{app}

\section{Objetivo de la aplicación}

Con el fin de facilitar el análisis e interpretación de los resultados obtenidos en los experimentos, se ha desarrollado una aplicación web interactiva que permite visualizar de forma dinámica las predicciones generadas por los modelos.

La aplicación no realiza entrenamiento en tiempo real, sino que consume los archivos generados automáticamente por el pipeline de modelado (ficheros CSV ubicados en el directorio \texttt{reports/}). De este modo, se mantiene una clara separación entre la fase de entrenamiento y la fase de visualización.

Esta arquitectura garantiza reproducibilidad, modularidad y simplicidad en el mantenimiento del sistema.

\section{Arquitectura del sistema}

La aplicación se integra como la capa final del flujo de trabajo desarrollado a lo largo del proyecto. La arquitectura general puede resumirse de la siguiente manera:

\begin{enumerate}
    \item Descarga de datos mediante la API de Eurostat.
    \item Transformación y cálculo de tasas de variación trimestral.
    \item Construcción del dataset unificado con retardos.
    \item Entrenamiento de modelos predictivos.
    \item Generación de métricas y predicciones en formato CSV.
    \item Visualización interactiva mediante la aplicación web.
\end{enumerate}

Esta estructura modular permite escalar el análisis a nuevos países o modelos sin modificar la lógica de la aplicación.

\section{Interfaz y funcionalidades}

La barra lateral de la aplicación expone tres selectores encadenados que determinan
completamente la vista mostrada:

\begin{itemize}
    \item \textbf{Bloque experimental}: permite elegir entre los tres bloques de
          experimentos implementados --- Bloque~0 (modelos por país), Bloque~1 (panel
          combinado) y Bloque~2 (panel con variable \emph{geo} como feature). La selección
          determina qué conjunto de ficheros CSV se carga.
    \item \textbf{País}: España (ES), Francia (FR) o Italia (IT). Filtra tanto las
          predicciones como las métricas al país seleccionado.
    \item \textbf{Modelo}: configuración base o ampliada. El nombre de la opción se
          adapta al bloque activo (p.ej. ``XGBoost base'' en el Bloque~0,
          ``Panel combinado ampliado'' en el Bloque~1).
\end{itemize}

A partir de estos tres selectores, la aplicación construye dinámicamente las rutas a los
ficheros CSV correspondientes, siguiendo la convención de nombres recogida en la
Tabla~\ref{tab:app_csvs}.

\begin{table}[H]
\centering
\caption{Ficheros CSV leídos por la aplicación según el bloque activo}
\label{tab:app_csvs}
\small
\begin{tabular}{lll}
\hline
\textbf{Bloque} & \textbf{Predicciones} & \textbf{Métricas} \\
\hline
0 & \texttt{predictions\_xgboost\_\{base|ext\}\_\{geo\}.csv}       & \texttt{metrics\_xgboost\_\{geo\}\_compare.csv} \\
1 & \texttt{predictions\_panel\_combined\_\{base|ext\}\_\{geo\}.csv} & \texttt{metrics\_panel\_combined\_compare.csv}  \\
2 & \texttt{predictions\_panel\_geo\_\{base|ext\}\_\{geo\}.csv}      & \texttt{metrics\_panel\_geo\_compare.csv}       \\
\hline
\end{tabular}
\end{table}

En los Bloques~1 y~2 el fichero de métricas contiene los resultados de los tres países en
un único CSV (columna \texttt{geo}); la aplicación filtra automáticamente las filas
correspondientes al país seleccionado. En el Bloque~0 cada país tiene su propio fichero
de métricas.

Una vez cargados los datos, el área principal muestra:

\begin{itemize}
    \item \textbf{Indicadores KPI}: modelo activo, MAE y RMSE en puntos porcentuales,
          y número de variables del modelo.
    \item \textbf{Gráfica de predicción}: comparación entre el PIB QoQ real y el
          predicho a lo largo del conjunto de prueba, generada con Plotly y con
          interactividad nativa (zoom, hover por trimestre).
    \item \textbf{Gráfica de errores}: evolución del error de predicción trimestre a
          trimestre con línea de referencia en cero.
    \item \textbf{Tablas opcionales}: predicciones completas y métricas comparativas,
          activables mediante checkbox en la barra lateral.
\end{itemize}

\section{Decisiones técnicas}

La aplicación ha sido desarrollada con Streamlit, seleccionado por su integración
directa con el ecosistema Python del proyecto y su capacidad para construir interfaces
interactivas sin necesidad de código JavaScript o HTML adicional.

Las visualizaciones se implementan con Plotly, que proporciona interactividad nativa
(zoom, selección de rango, inspección de valores por trimestre) sin coste de desarrollo
adicional respecto a bibliotecas estáticas como Matplotlib.

Se ha optado por una arquitectura desacoplada en la que la aplicación únicamente lee
los CSV generados por los scripts de entrenamiento. Esta separación tiene dos ventajas
concretas: permite actualizar los resultados simplemente re-ejecutando el script de
entrenamiento correspondiente sin tocar la aplicación, y facilita la incorporación de
nuevos bloques experimentales añadiendo únicamente los CSV con la convención de nombres
establecida.
